% ========================================================================= %
% NOWOCZESNY SZABLON CV W SYSTEMIE LATEX (DWUKOLUMNOWY)
% Wersja zoptymalizowana pod systemy ATS (Applicant Tracking Systems).
%
% Instrukcja użycia:
% 1. Podmień dane w sekcjach [PLACEHOLDER] na własne.
% 2. Zdjęcie: Wgraj plik (np. photo_01.png) do folderu projektu.
% 3. Kompilacja: Użyj silnika pdfLaTeX (np. w serwisie Overleaf).
% ========================================================================= %

\documentclass[10pt,a4paper,sans]{article}

% --- Pakiety podstawowe ---
\usepackage[utf8]{inputenc}
\usepackage[T1]{fontenc}
\usepackage[polish]{babel}
\usepackage{geometry}
\usepackage{xcolor}
\usepackage{enumitem}
\usepackage{fontawesome5}  % Ikonki (telefon, github, email)
\usepackage{graphicx}
\usepackage{hyperref}
\usepackage{titlesec}
\usepackage{tikz}

% --- Konfiguracja strony ---
\geometry{left=0cm, top=0cm, right=0cm, bottom=0cm} % Marginesy sterowane ręcznie przez minipage
\setlength{\parindent}{0pt} % Wyłączenie wcięć akapitowych
\pagestyle{empty} % Wyłączenie numeracji stron

% --- Definicje kolorów ---
\definecolor{sidebarcolor}{HTML}{F2F2F2} % Jasnoszary kolor lewego paska
\definecolor{maindark}{HTML}{2E2E2E}     % Ciemnoszary kolor dla nagłówków i tekstu

% --- Konfiguracja nagłówków sekcji ---
\titleformat{\section}{\large\bfseries}{}{0pt}{\MakeUppercase}[\titlerule]
\titlespacing{\section}{0pt}{14pt}{8pt}

% --- Konfiguracja linków ---
\hypersetup{
    colorlinks=true,
    urlcolor=black,
    linkcolor=black
}

% ========================================================================= %
% --- POCZĄTEK DOKUMENTU ---
% ========================================================================= %
\begin{document}%
%
% --- Tło lewej kolumny (Sidebar) ---
\begin{tikzpicture}[remember picture, overlay]
    \fill[sidebarcolor] (current page.north west) rectangle ([xshift=6cm]current page.south west);
\end{tikzpicture}%
%
% Kluczowe: brak pustych linii między tikz a minipage zapobiega pustym stronom.
\noindent
\begin{minipage}[t]{6cm} % --- LEWA KOLUMNA (SZEROKOŚĆ 6CM) ---
    \vspace{1cm}
    \centering

    % --- SEKCJA: ZDJĘCIE ---
    \begin{tikzpicture}
        \clip (0,0) circle (2cm); % Przycięcie zdjęcia do koła
        \node at (0,0) {\includegraphics[width=4cm]{photo_01.png}}; % Podmień photo_01.png na swój plik
        % Jeśli nie chcesz zdjęcia, zakomentuj linię powyżej i odkomentuj poniższą:
        % \draw[fill=white] (0,0) circle (2cm) node {FOTO};
    \end{tikzpicture}

    \vspace{1cm}

    \raggedright
    \hspace{0.5cm}
    \begin{minipage}{5cm} % Kontener na tekst wewnątrz kolumny
        \section*{Kontakt}
        \small
        \faPhone \hspace{0.2cm} +48 000 000 000 \\
        \faEnvelope \hspace{0.2cm} \href{mailto:email@example.com}{email@example.com} \\
        \faGithub \hspace{0.2cm} \href{https://github.com/YourUsername}{github.com/YourUsername} \\
        \faMapMarker \hspace{0.2cm} Miasto, Polska

        \section*{Technologie}
        \textbf{Języki:} \\
        $\bullet$ Java \\
        $\bullet$ C\# (Unity) \\
        $\bullet$ SQL \\
        $\bullet$ C++ \\
        $\bullet$ Python \\
        $\bullet$ GDScript (Godot) \\

        \vspace{2pt}
        \textbf{Narzędzia \& Systemy:} \\
        $\bullet$ Git (GitHub/GitLab) \\
        $\bullet$ Unity \\
        $\bullet$ SFML \\
        $\bullet$ Docker (podstawy) \\
        $\bullet$ Linux (Fedora), Windows \\
        $\bullet$ IntelliJ IDEA

        \section*{Języki}
        $\bullet$ Angielski (B2) \\
        $\bullet$ Polski (Ojczysty)

        \section*{Zainteresowania}
        $\bullet$ Gamedev i światotwórstwo \\
        $\bullet$ Grafika 2D (retro) i modelowanie 3D \\
        $\bullet$ Hardware i systemy (Linux) \\
        $\bullet$ Sport
    \end{minipage}
\end{minipage}%
\hspace{0.5cm}%
\begin{minipage}[t]{13.5cm} % --- PRAWA KOLUMNA (GŁÓWNA) ---
    \vspace{1.5cm}

    % --- NAGŁÓWEK: IMIĘ I NAZWISKO ---
    {\Huge \textbf{Kamil Śliwa}} \\
    {\large \textcolor{maindark}{Student Informatyki / Młodszy Programista}}

    \vspace{0.5cm}

    \section*{Podsumowanie zawodowe}
    Student 3. roku informatyki z pierwszym doświadczeniem komercyjnym zdobytym podczas praktyk programistycznych w obszarze gamedev. Specjalizuję się w programowaniu w językach Java oraz C\# (Unity), posiadam doświadczenie w pracy zespołowej oraz rozwijaniu i utrzymywaniu istniejących projektów. Realizowałem projekty webowe i gry, obejmujące implementację logiki aplikacji, mechanik rozgrywki oraz pracę z systemem kontroli wersji Git. Szukam możliwości rozwoju jako Junior Developer (Java / Unity).

    \section*{Doświadczenie}
    \textbf{[Nazwa Firmy / Praktyk]} \hfill Lipiec 2025 -- Sierpień 2025 \\
    \textit{Programista Praktykant} \\
    $\bullet$ Implementacja i rozwój mechanik rozgrywki w silniku Unity (C\#). \\
    $\bullet$ Analiza, diagnozowanie oraz naprawa błędów zgłaszanych w projektach (bug fixing). \\
    $\bullet$ Współpraca w zespole programistycznym z wykorzystaniem GitLab (version control, code review). \\
    $\bullet$ Utrzymanie i rozwój istniejącego kodu w środowisku produkcyjnym.

    \section*{Edukacja}
    \textbf{[Nazwa Uczelni / Uniwersytetu]} \hfill 2023 -- obecnie \\
    \textit{Kierunek: Informatyka (Studia inżynierskie)} \\

    \vspace{10pt}
    \textbf{[Nazwa Szkoły Średniej / Technikum]} \hfill 2019 -- 2023 \\
    \textit{Tytuł: Technik Informatyk} \\
    $\bullet$ Uzyskanie certyfikatów kwalifikacji zawodowych (INF.02, INF.03).

    \section*{Wybrane Projekty}
    \textbf{[Nazwa Projektu 1] (Aplikacja Webowa)} \\
    $\bullet$ Aplikacja do zarządzania i podziału wydatków między użytkownikami. \\
    $\bullet$ Implementacja logiki biznesowej oraz operacji na danych (CRUD). \\
    $\bullet$ Projekt obejmował pracę z bazą danych oraz obsługę użytkowników.

    \vspace{5pt}
    \textbf{[Nazwa Projektu 2] (Aplikacja Webowa)} \\
    $\bullet$ Projekt symulujący system losowania przedmiotów (lootbox) oraz zarządzania zasobami użytkownika. \\
    $\bullet$ Implementacja mechaniki losowania oraz systemu wymiany wirtualnych przedmiotów.

    \vspace{5pt}
    \textbf{Autorskie Projekty Gamedevowe} \\
    $\bullet$ Tworzenie projektów gier od podstaw: projektowanie, implementacja oraz podstawowa oprawa graficzna (2D/3D). \\
    $\bullet$ Implementacja klonu gry Wolfenstein 3D (mechaniki retro, rendering, logika ruchu). \\
    $\bullet$ Stworzenie gry typu Roguelike w języku Java (logika gry, system walki, zarządzanie stanem).

    \section*{Umiejętności Dodatkowe}
    \begin{itemize}[leftmargin=*, noitemsep]
        \item Doświadczenie w pracy zespołowej oraz korzystaniu z systemów kontroli wersji (Git).
        \item Umiejętność analizy problemów oraz debugowania kodu.
        \item Dobra organizacja pracy własnej i samodzielność w realizacji zadań.
        \item Szybka adaptacja do nowych technologii i środowisk programistycznych.
    \end{itemize}

    \vfill % Wypchnięcie klauzuli na dół strony

    % --- STOPKA: KLAUZULA RODO ---
    \begin{center}
        \scriptsize \textcolor{gray}{
        Wyrażam zgodę na przetwarzanie moich danych osobowych dla potrzeb niezbędnych do realizacji procesu tej oraz przyszłych rekrutacji (zgodnie z ustawą z dnia 10 maja 2018 roku o ochronie danych osobowych (Dz. Ustaw z 2018, poz. 1000) oraz zgodnie z Rozporządzeniem Parlamentu Europejskiego i Rady (UE) 2016/679 z dnia 27 kwietnia 2016 r. w sprawie ochrony osób fizycznych w związku z przetwarzaniem danych osobowych i w sprawie swobodnego przepływu takich danych oraz uchylenia dyrektywy 95/46/WE (RODO)).
        }
    \end{center}

    \vspace{0.5cm}
\end{minipage}

\end{document}